% =============================================================
%  Chapter -- Link
%  Subsections: Operating point and direction, DC circuit, Solution
%  procedures, Link totals and checks, Worked example (Table points).
%  Tables and values: fig/points_table.tex, fig/link_values.tex, generated
%  by fig/link_example.py from tools/example_data.py.
%  Sources: [1] hvdc_lf, src/hvdc_lf/link/case.py, link/solver.py.
% =============================================================
\section{Link}
\label{sec:link}

A link consists of two stations A and B, each calculated as in Section~\ref{sec:station}, and the DC line of
Section~\ref{sec:dc-line} between their DC terminals. The link configuration holds the two stations, the DC line and
the topology; the rectifier and the inverter are not part of it. Both stations may have the same design or different
ones.

\subsection{Operating Point and Direction}
\label{sec:link-point}

The operating point of a link gives, at both PCCs, the voltage, frequency, reactive power, tap position and auxiliary
power, the DC voltage $U_{dpe\_rec}$ at the rectifier DC terminal, the conductor temperature and the cooling data of
both converters. The active power is given at exactly one of four points, signed as for a single station:
$P_{pcc\_a}$, $P_{pcc\_b}$ (PCC of station A or B) or $P_{dc\_a}$, $P_{dc\_b}$ (DC terminal of station A or B, all
poles). The station at which the power is given is the rectifier if the power is negative, the inverter if it is
positive; the other station takes the other role. Seen from the power flow, the given point is one of four
references, which select the solution procedure of Section~\ref{sec:link-procedures}, Table~\ref{tab:link-reference}.

\begin{table}[!ht]
  \centering
  \small
  \begin{tabular}{@{}l l l l@{}}
    \toprule
    \textbf{Given} & \multicolumn{1}{c}{\textbf{Sign}} & \multicolumn{1}{c}{\textbf{Rectifier}}
      & \multicolumn{1}{c}{\textbf{Reference}} \\
    \midrule
    $P_{pcc\_a}$ & $< 0$ & A & sending AC \\
    $P_{pcc\_a}$ & $> 0$ & B & receiving AC \\
    $P_{pcc\_b}$ & $< 0$ & B & sending AC \\
    $P_{pcc\_b}$ & $> 0$ & A & receiving AC \\
    $P_{dc\_a}$  & $< 0$ & A & sending DC \\
    $P_{dc\_a}$  & $> 0$ & B & receiving DC \\
    $P_{dc\_b}$  & $< 0$ & B & sending DC \\
    $P_{dc\_b}$  & $> 0$ & A & receiving DC \\
    \bottomrule
  \end{tabular}
  \caption{Rectifier and power reference from the point and the sign of the given active power}
  \label{tab:link-reference}
\end{table}

\subsection{DC Circuit}
\label{sec:link-dc}

The DC circuit of one pole runs from the converter terminals of the rectifier through its DC reactor, the DC line
and the DC reactor of the inverter to the converter terminals of the inverter. For the current $I_d$ and the given
$U_{dpe\_rec}$,
\begin{align}
  U_{dpe\_conv\_rec} &= U_{dpe\_rec} + R_{dcr\_rec} \cdot I_d , \\
  U_{dpe\_inv} &= U_{dpe\_rec} - R_{pole} \cdot I_d , \\
  U_{dpe\_conv\_inv} &= U_{dpe\_inv} - R_{dcr\_inv} \cdot I_d ,
\end{align}
with $U_{dpe\_inv}$ the voltage at the inverter DC terminal (line end). The converters see
$U_{dc} = k \cdot U_{dpe\_conv}$ and carry $|P_{dc\_cv}| = k \cdot U_{dpe\_conv} \cdot I_d$ each.

\subsection{Solution Procedures}
\label{sec:link-procedures}

The four procedures use the station chain and the secant search on $P_{pcc}$ of
Section~\ref{sec:station-procedures} for a station whose converter DC power is known. The link secants have the
tolerance $10^{-6}$\,MW and at most 50 iterations.

\paragraph{Sending AC.} $P_{pcc}$ of the rectifier is given.
\begin{enumerate}
  \item Rectifier chain with the given $P_{pcc}$ and the fixed point on $U_{dpe\_conv\_rec}$ of
        Section~\ref{sec:station-procedures}, giving $I_d$.
  \item DC circuit, giving $U_{dpe\_conv\_inv}$.
  \item Inverter: secant on its $P_{pcc}$ until its converter DC power is $k \cdot U_{dpe\_conv\_inv} \cdot I_d$.
\end{enumerate}

\paragraph{Sending DC.} $P_{dc}$ at the rectifier DC terminal is given.
\begin{enumerate}
  \item $I_d = |P_{dc}| / (n_{conv} \cdot k \cdot U_{dpe\_rec})$ and the DC circuit, giving both converter voltages.
  \item Rectifier: secant on its $P_{pcc}$ until its converter DC power is $-k \cdot U_{dpe\_conv\_rec} \cdot I_d$.
  \item Inverter: secant as for sending AC, independent of step 2.
\end{enumerate}

\paragraph{Receiving DC.} $P_{dc}$ at the inverter DC terminal is given.
\begin{enumerate}
  \item Secant on $I_d$ until $n_{conv} \cdot k \cdot U_{dpe\_inv}(I_d) \cdot I_d = P_{dc}$, starting from
        $I_d = P_{dc} / (n_{conv} \cdot k \cdot U_{dpe\_rec})$ and 1.01 times that value. This is the receiving end
        current of Section~\ref{sec:dc-line}, solved numerically through the DC circuit.
  \item Both stations as for sending DC.
\end{enumerate}

\paragraph{Receiving AC.} $P_{pcc}$ of the inverter is given.
\begin{enumerate}
  \item Inverter chain with the given $P_{pcc}$ at the converter voltage $U_{dpe\_conv\_inv}$, starting with
        $U_{dpe\_conv\_inv} = U_{dpe\_rec}$, giving $I_d$.
  \item DC circuit with this $I_d$, giving a new $U_{dpe\_conv\_inv}$; repeat from step 1 until the change is below
        $10^{-6}$\,kV.
  \item Rectifier: secant as for sending DC.
\end{enumerate}

\subsection{Link Totals and Checks}
\label{sec:link-totals}

The powers at the DC terminals are the converter DC powers of all converters plus the DC reactor losses of all pole
conductors, both in the sign convention of Section~\ref{sec:conventions}. The link losses and the efficiency are
\begin{equation}
  P_{loss\_link} = -\left( P_{pcc\_rec} + P_{pcc\_inv} \right) , \qquad
  \eta = \frac{P_{pcc\_inv}}{-P_{pcc\_rec}} .
\end{equation}
The losses must equal the sum of the station losses of both stations and the losses of the $n_{pc}$ pole conductors;
a difference above $10^{-3}$\,MW is a warning. The limit warnings of both stations (Section~\ref{sec:station}) are
collected in the link result.

\subsection{Worked Example}
\label{sec:link-example}

Two stations of the design of Table~\ref{tab:example-data} are joined by the cable of
Section~\ref{sec:dc-line-example}, with the station operating point of Section~\ref{sec:example-data} at both PCCs and
$U_{dpe\_rec} = 400$\,kV. With $P_{pcc\_a} = -1000$\,MW station A is the rectifier; $I_d = \lkId$\,kA, the inverter DC
terminal voltage is $\lkUinv$\,kV, the line losses are $\lkLine$\,MW and the link losses $\lkLoss$\,MW, at an
efficiency of $\lkEta$\,\%. Table~\ref{tab:link-points} gives this operating point with the power given at each of
the four points, rows 1 to 4, and two further points. Rows 1 to 4 give the same operating point, with the power taken
from row 1. In row 5 the power at the PCC of station A is positive, so station B is the rectifier; the DC voltage of
400\,kV is then given at station B, and the current and the losses are higher than in row 1 for the same power at
PCC A, because 1000\,MW are now delivered to grid A instead of taken from it. Row 6 is row 1 mirrored: station B takes
1000\,MW from its grid.

\begin{table}[!ht]
  \centering
  \footnotesize
  \setlength{\tabcolsep}{4pt}
  \begin{tabular}{@{}l c r r r r r r r@{}}
    \toprule
    \textbf{Power given} & \multicolumn{1}{c}{\textbf{Rect.}} & \multicolumn{1}{c}{\textbf{$P_{pcc\_a}$}}
      & \multicolumn{1}{c}{\textbf{$P_{dc\_a}$}} & \multicolumn{1}{c}{\textbf{$P_{dc\_b}$}}
      & \multicolumn{1}{c}{\textbf{$P_{pcc\_b}$}} & \multicolumn{1}{c}{\textbf{$I_d$}}
      & \multicolumn{1}{c}{\textbf{Losses}} & \multicolumn{1}{c}{\textbf{Iter.}} \\
    & & \multicolumn{1}{c}{[MW]} & \multicolumn{1}{c}{[MW]} & \multicolumn{1}{c}{[MW]} & \multicolumn{1}{c}{[MW]}
      & \multicolumn{1}{c}{[kA]} & \multicolumn{1}{c}{[MW]} & \\
    \midrule
    \tablerows{link/fig/points_table}
    \bottomrule
  \end{tabular}
  \caption{Link operating points with the active power given at different points, $U_{dpe\_rec} = 400$\,kV,
           $U_{pcc} = 345$\,kV and $Q_{pcc} = 0$ at both stations}
  \label{tab:link-points}
\end{table}
\FloatBarrier
